\clearpage
\begin{examplebox}
\textbf{\textcolor{CExample}{例 1（最基础：两次递减，为什么最后递增？）}}
\begin{center}\begin{tikzpicture}[font=\small,>=stealth,line join=round]
\draw[gray!40,rounded corners=3pt] (0,0) rectangle (14.2,9.0);
\node at (7.1,8.4) {先减一次，再减一次：大小顺序恢复};
\draw[->,gray] (1.0,7.2)--(7.5,7.2) node[right]{$x$};
\fill[blue] (2.3200000000000003,7.2) circle(2.2pt);
\node[above left=7pt,fill=white,inner sep=1pt] at (2.3200000000000003,7.2) {$2$};
\draw[blue,dashed,thick,->,shorten <=3pt,shorten >=3pt] (2.3200000000000003,7.2)--(5.98,4.7);
\fill[orange!85!black] (4.15,7.2) circle(2.2pt);
\node[above right=7pt,fill=white,inner sep=1pt] at (4.15,7.2) {$5$};
\draw[orange!85!black,dashed,thick,->,shorten <=3pt,shorten >=3pt] (4.15,7.2)--(4.15,4.7);
\draw[->,gray] (1.0,4.7)--(7.5,4.7) node[right]{$t$};
\fill[blue] (5.98,4.7) circle(2.2pt);
\node[above right=7pt,fill=white,inner sep=1pt] at (5.98,4.7) {$8$};
\draw[blue,dashed,thick,->,shorten <=3pt,shorten >=3pt] (5.98,4.7)--(2.3200000000000003,2.2);
\fill[orange!85!black] (4.15,4.7) circle(2.2pt);
\node[above left=7pt,fill=white,inner sep=1pt] at (4.15,4.7) {$5$};
\draw[orange!85!black,dashed,thick,->,shorten <=3pt,shorten >=3pt] (4.15,4.7)--(4.15,2.2);
\draw[->,gray] (1.0,2.2)--(7.5,2.2) node[right]{$y$};
\fill[blue] (2.3200000000000003,2.2) circle(2.2pt);
\node[below left=7pt,fill=white,inner sep=1pt] at (2.3200000000000003,2.2) {$2$};
\fill[orange!85!black] (4.15,2.2) circle(2.2pt);
\node[below right=7pt,fill=white,inner sep=1pt] at (4.15,2.2) {$5$};
\node[align=left,anchor=west] at (8.4,4.7){$g(x)=10-x$：减\\[12pt]$f(t)=10-t$：减\\[14pt]$H(x)=x$：增\\[14pt]$2<5$\\[9pt]$8>5$\\[9pt]$2<5$};
\node[text width=12.8cm,align=center] at (7.1,.65) {蓝点与橙点各自代表同一个输入；各行数轴向右增大，连接线追踪实际数值。};
\end{tikzpicture}\end{center}



\textbf{题目：}判断 $H(x)=10-(10-x)$ 的单调性，不要先化简。

\textbf{解：}令 $g(x)=10-x$，$f(t)=10-t$。
当 $x$ 增大，$g(x)$ 减小，所以内层递减；当 $t$ 增大，$f(t)$ 减小，所以外层递减。
两次递减，整体递增。再化简检查：
\[
 H(x)=10-10+x=x,
\]
确实递增。\textbf{此例教会的是为什么“减减得增”，但不能代替一般证明。}

\medskip
\end{examplebox}
\clearpage
\begin{examplebox}
\textbf{\textcolor{CExample}{例 2（常见高考题：指数套二次函数）}}
\begin{center}\begin{tikzpicture}[font=\small,>=stealth,line join=round]
\draw[gray!40,rounded corners=3pt] (0,0) rectangle (14.2,11.3);
\node at (7.1,10.700000000000001) {例2：外层递减，内层最低点变成整体最高点};

\begin{scope}[shift={(2,7.3)},x=1.2cm,y=.55cm]
\draw[->,gray] (-1.5,0)--(3.6,0) node[right]{$x$};\draw[->,gray] (0,-1.6)--(0,4.5) node[above]{$t$};
\draw[blue,very thick,domain=-1:3,samples=100] plot(\x,{(\x-1)^2-1});\fill[orange!85!black] (1,-1) circle(2pt);\draw[gray,dotted] (1,-1)--(1,3.7);
\end{scope}
\begin{scope}[shift={(2,2.1)},x=1.2cm,y=1.0cm]
\draw[->,gray] (-1.5,0)--(3.6,0) node[right]{$x$};\draw[->,gray] (0,0)--(0,3.7) node[above]{$H$};
\draw[violet,very thick,domain=-1.3:3.3,samples=150] plot(\x,{pow(3,1-(\x-1)^2)});\fill[orange!85!black] (1,3) circle(2pt);\draw[gray,dotted] (1,0)--(1,3);\node[below=5pt] at (1,0){$1$};
\end{scope}
\node[align=left,anchor=west] at (8.3,8.1){$g(x)=(x-1)^2-1$\\[12pt]$x=1$：内层最小 $-1$\\[12pt]$g'(x)=2(x-1)$};
\node[align=left,anchor=west] at (8.3,3.9){$H(x)=3^{-g(x)}$\\[12pt]外层 $3^{-t}$ 递减\\[12pt]内层的谷，对应整体的峰\\[12pt]$H(1)=3$ 为最大值};

\node[text width=12.8cm,align=center] at (7.1,.65) {两图共享相同的 $x$ 刻度；纵轴分别为内层 $t$ 与整体函数值 $H$。};
\end{tikzpicture}\end{center}



\textbf{题目：}求 $H(x)=\left(\frac13\right)^{x^2-2x}$ 的单调区间，并求最大值。

\textbf{第一步：}底数为正且不等于 $1$，指数可以是任意实数，所以定义域为 $\mathbb R$。

\textbf{第二步：}令
\[
 t=x^2-2x=(x-1)^2-1,\qquad H(x)=\left(\frac13\right)^t.
\]
内层是一条开口向上的抛物线：在 $x\le1$ 这一侧递减，在 $x\ge1$ 这一侧递增。

\textbf{第三步：}外层 $\left(\frac13\right)^t$ 对所有实数 $t$ 都递减。
因此左边“减、减、增”，右边“增、减、减”。
\[
\boxed{H(x)\text{ 在 }(-\infty,1]\text{ 上递增，在 }[1,+\infty)\text{ 上递减。}}
\]

\textbf{第四步（亮点）：}内层在 $x=1$ 取最小值 $-1$，
外层递减，故整体在这里取\textbf{最大值}
\[
 H(1)=\left(\frac13\right)^{-1}=3.
\]
\textbf{顺带记住：}外层递减时，里面的最小值往往对应外面的最大值。

\medskip
\end{examplebox}
\clearpage
\begin{examplebox}
\textbf{\textcolor{CExample}{例 3（先查定义域：对数套二次函数）}}
\begin{center}\begin{tikzpicture}[font=\small,>=stealth,line join=round]
\draw[gray!40,rounded corners=3pt] (0,0) rectangle (14.2,9.7);
\node at (7.1,9.1) {例3：先删掉真数不为正的区间，再判断两支增减};

\begin{scope}[shift={(3.7,4.6)},x=.68cm,y=.85cm]
\fill[gray!10] (-1,-3) rectangle (3,3);
\draw[->,gray] (-3.2,0)--(5.6,0) node[right]{$x$};\draw[->,gray] (0,-3.3)--(0,3.4) node[above]{$H$};
\draw[dashed,gray] (-1,-3)--(-1,3) (3,-3)--(3,3);
\begin{scope}\clip (-3.2,-3.1) rectangle (5.4,3.1);\draw[blue,very thick,domain=-3:-1.001,samples=150] plot(\x,{ln(\x*\x-2*\x-3)});
\draw[blue,very thick,domain=3.001:5.3,samples=150] plot(\x,{ln(\x*\x-2*\x-3)});
\end{scope}\node[below=5pt] at (-1,-3){$-1$};\node[below=5pt] at (3,-3){$3$};
\end{scope}
\node[align=left,anchor=west] at (8.4,5){$H=\ln(x^2-2x-3)$\\[12pt]真数要求 $(x+1)(x-3)>0$\\[12pt]左支 $x<-1$：减\\[12pt]右支 $x>3$：增\\[12pt]浅色带内没有定义};

\node[text width=12.8cm,align=center] at (7.1,.65) {不能把中间没有定义的部分补进来，也不能把两支当作一个连续区间。};
\end{tikzpicture}\end{center}



\textbf{题目：}求 $H(x)=\ln(x^2-2x-3)$ 的单调区间。

\textbf{第一步：}对数里面必须大于 $0$：
\[
 x^2-2x-3=(x-3)(x+1)>0
 \quad\Longrightarrow\quad x<-1\ \text{或}\ x>3.
\]
原函数只在 $(-\infty,-1)\cup(3,+\infty)$ 有定义。

\textbf{第二步：}令 $t=(x-1)^2-4$。
内层以 $x=1$ 为分界：左边递减，右边递增。
但能用的区间只有前面求出的两段。

\textbf{第三步：}外层 $\ln t$ 在 $t>0$ 时递增。
因此在 $(-\infty,-1)$ 上“减、增、减”；
在 $(3,+\infty)$ 上“增、增、增”。
\[
\boxed{\text{递减区间 }(-\infty,-1),\qquad
       \text{递增区间 }(3,+\infty).}
\]
\textbf{提醒：}不要把中间没定义的一段也算进去，更不能把两侧无定义地连成一个大区间。

\medskip
\end{examplebox}
\clearpage
\begin{examplebox}
\textbf{\textcolor{CExample}{例 4（逆向应用：含参数时两个条件都不能漏）}}
\begin{center}\begin{tikzpicture}[font=\small,>=stealth,line join=round]
\draw[gray!40,rounded corners=3pt] (0,0) rectangle (14.2,9.8);
\node at (7.1,9.200000000000001) {例4：递增还不够，真数也必须始终大于零};

\begin{scope}[shift={(-2.2,2.9)},x=2.0cm,y=.48cm]
\draw[->,gray] (1.8,0)--(4.7,0) node[right]{$x$};\draw[->,gray] (2,-2)--(2,10.5) node[above]{真数};
\draw[gray,dashed] (2,-2)--(2,10);\node[below left=6pt] at (2,0){$2$};
\draw[blue,thick,domain=2:4.3,samples=100] plot(\x,{\x*\x-2*\x+1});
\draw[teal!80!black,very thick,domain=2:4.3,samples=100] plot(\x,{\x*\x-2.5*\x+1});
\draw[orange!85!black,dashed,thick,domain=2:4.3,samples=100] plot(\x,{\x*\x-3*\x+1});
\draw[blue,fill=white] (2,1) circle(2.4pt);\draw[teal!80!black,fill=white] (2,0) circle(2.4pt);\draw[orange!85!black,fill=white] (2,-1) circle(2.4pt);
\end{scope}
\node[align=left,anchor=west] at (8.1,5){$t=x^2-ax+1,\ x>2$\\[12pt]\textcolor{blue}{$a=2$：真数恒正}\\[12pt]\textcolor{teal!80!black}{$a=5/2$：端点趋于0，允许}\\[12pt]\textcolor{orange!85!black}{$a=3$：靠近2时真数为负}\\[14pt]最终 $a\le5/2$};

\node[text width=12.8cm,align=center] at (7.1,.65) {三条线均递增，只有蓝、绿线的真数恒正；空心点不取。};
\end{tikzpicture}\end{center}



\textbf{题目：}已知 $H(x)=\ln(x^2-ax+1)$ 在 $(2,+\infty)$ 上有定义，
且严格递增，求实数 $a$ 的范围。

\textbf{第一步：}设 $t=x^2-ax+1$。因为 $\ln t$ 在 $t>0$ 时严格递增，
要使整体在 $(2,+\infty)$ 上递增，里面的二次函数也要在这段上递增，
同时必须处处满足 $t>0$。

\textbf{第二步（递增条件）：}内层对称轴是 $x=\frac a2$。
开口向上，要使 $(2,+\infty)$ 位于其递增的一侧，需要
\[
 \frac a2\le2\quad\Longrightarrow\quad a\le4.
\]

\textbf{第三步（有定义条件）：}当内层在 $(2,+\infty)$ 递增时，
要使每个 $x>2$ 都有 $x^2-ax+1>0$，必须且只需紧靠左端点时不出现负值。
代入左端点的数值作检查：
\[
 2^2-2a+1=5-2a\ge0
 \quad\Longrightarrow\quad a\le\frac52.
\]
注意 $(2,+\infty)$ \textbf{不含} $x=2$，所以这里允许 $5-2a=0$。
此时由于二次函数随后严格增大，所有 $x>2$ 的真数仍然大于 $0$。

\textbf{第四步：}因为 $a\le\frac52$ 自动满足 $a\le4$，结果是
\[
 \boxed{a\in\left(-\infty,\frac52\right].}
\]
\textbf{验算边界：}当 $a=\frac52$，真数
$=(x-2)\left(x-\frac12\right)>0$（$x>2$），且随 $x$ 严格增大，故边界可以取到。

\medskip
\end{examplebox}
\clearpage
\begin{examplebox}
\textbf{\textcolor{CExample}{例 5（反例提醒：外层跨过掉头点，口诀会失效）}}
\begin{center}\begin{tikzpicture}[font=\small,>=stealth,line join=round]
\draw[gray!40,rounded corners=3pt] (0,0) rectangle (14.2,9.4);
\node at (7.1,8.8) {例5：内层跨过零，外层平方要分成两边判断};

\begin{scope}[shift={(2.5,2.3)},x=1.2cm,y=.85cm]
\draw[->,gray] (-1.5,0)--(3.8,0) node[right]{$x$};\draw[->,gray] (0,-.3)--(0,5.3) node[above]{$H$};
\draw[blue,very thick,domain=-1.2:3.2,samples=140] plot(\x,{(\x-1)^2});\draw[gray,dotted] (1,0)--(1,4.8);\fill[orange!85!black] (1,0) circle(2.4pt);\node[below=6pt] at (1,0){$1$};
\end{scope}
\node[align=left,anchor=west] at (8.3,5){$t=x-1$：内层一直增\\[12pt]$t<0$ 时平方递减\\[9pt]$t>0$ 时平方递增\\[14pt]$x\le1$：整体减\\[9pt]$x\ge1$：整体增};

\node[text width=12.8cm,align=center] at (7.1,.65) {外层转折发生在 $t=0$，换回原来的自变量位置就是 $x=1$。};
\end{tikzpicture}\end{center}



\textbf{题目：}判断 $H(x)=(x-1)^2$ 的增减性。

设 $g(x)=x-1$、$f(t)=t^2$。$g$ 在所有实数上递增；
可是外层 $f(t)=t^2$ 在 $t<0$ 时递减，在 $t>0$ 时递增。
因为 $t=x-1$，应当在 $x=1$ 处分开：
\[
\boxed{H(x)\text{ 在 }(-\infty,1]\text{ 上递减，在 }[1,+\infty)\text{ 上递增。}}
\]
\textbf{关键：}内层递增并不代表整体一定递增，得先看外层在\emph{内层实际输出}的范围内是增还是减。

\medskip
\end{examplebox}
\clearpage
\begin{examplebox}
\textbf{\textcolor{CExample}{例 6（多层推广：三个递减，最终怎样？）}}
\begin{center}\begin{tikzpicture}[font=\small,>=stealth,line join=round]
\draw[gray!40,rounded corners=3pt] (0,0) rectangle (14.2,11.6);
\node at (7.1,11.0) {例6：连续三次递减，最终大小关系相反};
\draw[->,gray] (1.0,9.6)--(7.5,9.6) node[right]{$x$};
\fill[blue] (2.3200000000000003,9.6) circle(2.2pt);
\node[above left=7pt,fill=white,inner sep=1pt] at (2.3200000000000003,9.6) {$2$};
\draw[blue,dashed,thick,->,shorten <=3pt,shorten >=3pt] (2.3200000000000003,9.6)--(5.98,7.0);
\fill[orange!85!black] (4.15,9.6) circle(2.2pt);
\node[above right=7pt,fill=white,inner sep=1pt] at (4.15,9.6) {$5$};
\draw[orange!85!black,dashed,thick,->,shorten <=3pt,shorten >=3pt] (4.15,9.6)--(4.15,7.0);
\draw[->,gray] (1.0,7.0)--(7.5,7.0) node[right]{$t$};
\fill[blue] (5.98,7.0) circle(2.2pt);
\node[above right=7pt,fill=white,inner sep=1pt] at (5.98,7.0) {$8$};
\draw[blue,dashed,thick,->,shorten <=3pt,shorten >=3pt] (5.98,7.0)--(2.3200000000000003,4.4);
\fill[orange!85!black] (4.15,7.0) circle(2.2pt);
\node[above left=7pt,fill=white,inner sep=1pt] at (4.15,7.0) {$5$};
\draw[orange!85!black,dashed,thick,->,shorten <=3pt,shorten >=3pt] (4.15,7.0)--(4.15,4.4);
\draw[->,gray] (1.0,4.4)--(7.5,4.4) node[right]{$u$};
\fill[blue] (2.3200000000000003,4.4) circle(2.2pt);
\node[above left=7pt,fill=white,inner sep=1pt] at (2.3200000000000003,4.4) {$2$};
\draw[blue,dashed,thick,->,shorten <=3pt,shorten >=3pt] (2.3200000000000003,4.4)--(5.98,1.8);
\fill[orange!85!black] (4.15,4.4) circle(2.2pt);
\node[above left=7pt,fill=white,inner sep=1pt] at (4.15,4.4) {$5$};
\draw[orange!85!black,dashed,thick,->,shorten <=3pt,shorten >=3pt] (4.15,4.4)--(4.15,1.8);
\draw[->,gray] (1.0,1.8)--(7.5,1.8) node[right]{$y$};
\fill[blue] (5.98,1.8) circle(2.2pt);
\node[below left=7pt,fill=white,inner sep=1pt] at (5.98,1.8) {$8$};
\fill[orange!85!black] (4.15,1.8) circle(2.2pt);
\node[below right=7pt,fill=white,inner sep=1pt] at (4.15,1.8) {$5$};
\node[align=left,anchor=west] at (8.4,6.0){每层都是 $10-$ 上一层\\[14pt]第一层：$2<5\to8>5$\\[12pt]第二层：$8>5\to2<5$\\[12pt]第三层：$2<5\to8>5$\\[14pt]三次反向，整体递减};
\node[text width=12.8cm,align=center] at (7.1,.65) {蓝点与橙点各自代表同一个输入；各行数轴向右增大，连接线追踪实际数值。};
\end{tikzpicture}\end{center}



\textbf{题目：}不展开式子，判断 $H(x)=10-\bigl[10-(10-x)\bigr]$ 的单调性。

从里向外看：$10-x$ 递减；$10-t$ 递减；最外面的 $10-u$ 也递减。
一共反转三次（奇数次），所以\textbf{整体递减}。
展开检查：
\[
 10-\bigl[10-(10-x)\bigr]=10-x,
\]
答案一致。\textbf{多层的快速技巧：只数递减出现了几层，前提是每层所在范围都单调。}
\end{examplebox}
